\documentclass[oneside,draft]{amsart}
\usepackage{amsmath,amssymb,amsthm,stmaryrd}
\usepackage[final]{hyperref}
\usepackage[final]{graphicx}
\usepackage{xcolor}
\usepackage{fixme}
\fxsetup{draft}
\fxusetheme{color}

\newcommand{\setcond}[2]{\left\{#1\mid#2\right\}}
\newcommand{\intrs}{\not\asymp}
\newcommand{\nintrs}{\asymp}
\newcommand{\funcoids}{\mathsf{FCD}}
\newcommand{\atoms}{\operatorname{atoms}}
\newcommand{\norm}[1]{\lVert #1\rVert}
\newcommand{\supfun}[1]{\left\langle#1\right\rangle}
\newcommand{\suprel}[1]{\left[#1\right]}
% \newcommand{\del}[1]{\textcolor{gray}{#1}}
\newenvironment{del}{\color{gray}}{}

\newtheorem{thm}{Theorem}
\newtheorem{lem}{Lemma}
\newtheorem{prop}{Proposition}
\newtheorem{note}{Note}
\newtheorem{obvious}{Obvious}
\newtheorem{cor}{Corollary}
\newtheorem{claim}{Claim}
\newtheorem{rem}{Remark}
\newtheorem{defn}{Definition}
\newtheorem{exer}{Exercise}

\title{[Second Attempt] A Proof of Kakeya's Conjecture}
\author{Victor Porton, ORCID 0000-0001-7064-7975}
\email{porton.victor@gmail.com}
\subjclass[2020]{28D99, 51M05, 54J99, 57N99, 51M15, 51M99}
\keywords{Kakeya conjecture, funcoids, dimensionality, diffeomorphism, geometry, general topology, long-standing open problem}

\begin{document}

\begin{abstract}
A rather short proof of Kakeya conjecture in arbitrary Euclidean spaces is presented. The proof uses several space transformations/projections. The conjecture is reduced to a Ka\-ke\-ya-li\-ke theorem on a hypersphere.
\end{abstract}

\maketitle

In this article, I prove Kakeya conjecture in an arbitrary Euclidean space.

\textbf{Kakeya set conjecture:} \cite{kakeya-long,tao-blog-kakeya} Define a \emph{Besicovitch} set in $\mathbb{R}^n$ to be a set which contains a unit line segment in every direction. Is it true that such sets necessarily have Hausdorff dimension and Minkowski dimension equal to~$n$?

We will prove only for Hausdorff dimension, because Minkowski dimension cannot be lower and having a proof for Hausdorff dimension also proves it for Minkowski one. $\dim$ operator will mean Hausdorff dimension.

\section{Kakeya proof}

Is it enough to prove it for lines instead of line segments? If for lines we prove that a finite part of the set has dimension~$n$, then for unit segments it easily?? follows.

So, we will solve for lines.

We can assume that all ``Besicovitch lines'' intersect a certain hyperplane (I call it the ``ground plane'').

Projections of lines to the ground plane form Besicovitch set of dimension~$n-1$. Assume it for induction.

??What's about the set of intersections of our lines with the ground plane?

On the ground plane, there is a set of dimensionality~$n-1$. Each point of that set lies on a line of our set of lines (in $n-1$, possibly not in $n$).

\section{``Political'' context}

I discovered ordered semigroup/semicategory actions in 2019. The world didn't react obeying like a sheep to decision of Russian Orthodoxes to kick people like me from normal life such as academic carrier. That's scary: the most valuable component is a missing component. Missing ordered semigroup actions amount to like half of world economy missing~\cite{osa-important}.

I don't know how Kakeya conjecture is used in the rest of math, but I use it to build a bridge between two about halves: academic math and the second lost half, through my glorification.

The advice: ``Solve some famous open problem to confirm utility of your theory.'' sounded for me like a mocking. But in the new reality, I indeed did it.

\bibliographystyle{amsplain}
\bibliography{kakeya}

\end{document}